\documentclass[runningheads]{llncs}

\usepackage[T1]{fontenc}
\usepackage[utf8]{inputenc}
\usepackage{graphicx}
\usepackage{booktabs}
\usepackage{tabularx}
\usepackage{url}
\usepackage{hyperref}
\hypersetup{colorlinks=true,linkcolor=blue,citecolor=blue,urlcolor=blue}

\title{DepChain: Dependable Permissioned Blockchain\\Stage 2 Report (Group 39)}

\author{Guilherme Ribeiro \and Tomás Morais}
\institute{Instituto Superior Técnico, University of Lisbon, Portugal\\
IST student nos.\ 103562, 103619}

\begin{document}
\maketitle

\begin{abstract}
DepChain is a simplified permissioned blockchain with Byzantine fault tolerance.
Stage~1 focused on ordering client requests with Basic HotStuff~\cite{hotstuff} over authenticated links.
Stage~2 refined the application semantics to support native cryptocurrency transfers (DepCoin),
smart contract execution on a Besu EVM, an ERC-20 token contract (IST Coin) hardened against approval frontrunning,
an Ethereum-like gas/fee model, block batching with a per-block gas cap, and JSON persistence of account state plus
IST Coin contract storage slots. We strengthened the client interface and tests to tolerate Byzantine clients and Byzantine replies,
while avoiding heavy full-chain sync for reads.
\keywords{Byzantine fault tolerance \and HotStuff \and blockchain \and EVM \and gas accounting}
\end{abstract}

\section{Introduction}
DepChain is a closed-membership blockchain service aiming for strong dependability guarantees in the presence of Byzantine faults.
The membership (replicas and their public keys) is static and known before start, and clients may be Byzantine but cannot compromise honest keys.
Stage~1 provided a BFT ordering layer based on Basic HotStuff and a secure client-to-replica request path.
Stage~2 extends the application layer to support transactions, fees, and smart contracts as required by the Stage~2 specification.

\section{System Overview}
\subsection{Stage 1 recap (consensus + links)}
DepChain uses a layered communication stack over UDP: a best-effort layer with retries (Fair Loss) and an Authenticated Perfect Link (APL).
The default APL mode attaches an RSA signature over the logical payload (membership publishes RSA verification keys).
As an alternative aligned with symmetric link authentication, replicas may derive per-neighbour secrets with ECDH on \texttt{secp256r1} (P-256)
and authenticate frames with HMAC-SHA256 instead; our multi-process key file appends the EC key material needed for that mode.
Both choices satisfy the design goal of forgery-resistant links among a closed membership; the report documents the option as a deliberate
trade-off (asymmetric broadcasts vs.\ compact MAC tags) rather than as a correctness requirement unique to one primitive.
Consensus is Basic HotStuff~\cite{hotstuff} with explicit \texttt{DECIDE} to define a single decision point, and quorum certificates (QC) built from $2f+1$ signatures.

\subsection{Stage 2 architecture additions}
Stage~2 introduces a transaction-aware blockchain ledger:
(i) a genesis block that initializes accounts and deploys an ERC-20~\cite{erc20} contract (IST Coin),
(ii) a world state tracking account balances and nonces,
(iii) a transaction executor that applies native transfers and contract deploy/calls using Hyperledger Besu's EVM,
and (iv) persistent storage for blocks using the same JSON spirit as genesis: each \texttt{LedgerBlock} includes hashes, executed
transactions, account snapshots, deployed runtime bytecode, and (for IST Coin) relevant contract storage slots.

\section{Transaction Model and Execution}
\subsection{Account model}
We follow the Ethereum account model with Externally Owned Accounts (EOA) and Contract Accounts.
Each account has an address, DepCoin balance, and nonce. Contract accounts additionally hold EVM bytecode and contract storage (key-value store).

\subsection{Transaction object}
Transactions include sender (\texttt{from}), recipient (\texttt{to}, or null for deploy), \texttt{nonce}, \texttt{value} (native transfer),
\texttt{gasPrice}, \texttt{gasLimit}, and \texttt{data} for contract deploy/call.
Transactions are signed by clients and verified by replicas before being admitted to the verified pool, preventing Byzantine clients from forging
operations on other accounts.

\subsection{Gas and fee model}
Each transaction is charged a fee deducted from the sender's DepCoin balance:
\[
\mathrm{fee} = \min(\mathrm{gasPrice}\cdot \mathrm{gasLimit},\ \mathrm{gasPrice}\cdot \mathrm{gasUsed})
\]
with \texttt{gasUsed} measured by the EVM for contract deploy/calls and set to an intrinsic constant for native transfers.
If \texttt{gasUsed} exceeds \texttt{gasLimit}, the transaction is aborted but the gas spent is not refunded (the cap still applies).
We guard against arithmetic overflow when computing $\mathrm{gasPrice}\cdot\mathrm{gasLimit}$.

\subsection{Smart contract execution (Besu EVM)}
Contract deployment and contract calls are executed using Hyperledger Besu~\cite{besu}.
For a known IST Coin creation bytecode, replicas install the compiled runtime bytecode and seed initial balances.
Contract addresses are derived deterministically from \texttt{(from, nonce)} using a stable project-specific function so that all replicas compute
the same contract address.

\section{Block Construction, Ordering, and Persistence}
\subsection{Batching and fee-based ordering}
For efficiency, leaders propose blocks containing multiple transactions. Transactions are ordered by \emph{fee offer},
placing higher-fee transactions earlier in the block. Ties are broken to preserve per-sender nonce ordering to avoid proposing invalid sequences.

\subsection{Per-block gas cap}
To avoid oversized blocks and approximate an Ethereum-like block ``weight'', we enforce a maximum gas budget per block.
The leader greedily selects transactions in priority order until adding the next transaction would exceed the gas cap, leaving the remainder in the mempool.

\subsection{Genesis and JSON persistence}
The system starts with a genesis JSON file that contains the initial world state (accounts, balances, nonces) and initial transactions (including
the IST Coin deployment). An optional \texttt{contracts} array allows bootstrapping arbitrary contract accounts by supplying runtime bytecode
and an explicit map of 32-byte storage slots (hex keys and values), without requiring a creation transaction in \texttt{transactions}.
This generalises ``code + storage at genesis'' beyond the IST-only bootstrap path. Each subsequent block points to its predecessor and is written as JSON (\texttt{BlockJsonStore}), including:
\begin{itemize}
  \item \texttt{state}: balances and nonces for all known accounts at the block head;
  \item \texttt{contractRuntimeHex}: contract address $\rightarrow$ runtime bytecode (hex) so replicas can reinstall code after load;
  \item \texttt{contractStorageHex}: for supported contracts (currently IST Coin only), contract address $\rightarrow$ map of 32-byte storage slots to 32-byte values
  (both as 64-character hex, no \texttt{0x} prefix).
\end{itemize}

\paragraph{IST Coin storage snapshot (pragmatic completeness).}
Full Besu \texttt{SimpleWorld} serialization would be large. We persist IST Coin storage by enumerating only the ERC-20 mapping slots used in our Solidity
(\texttt{\_balances} at slot~0 and \texttt{\_allowances} at slot~1) for a bounded set of known addresses (e.g., up to 32). Only non-zero slots are stored
(implicit zero otherwise). After load, replicas reinstall code from \texttt{contractRuntimeHex} and apply stored slots from \texttt{contractStorageHex}.

\paragraph{Relation to ``world state after the block'' (Stage~2 PDF).}
This satisfies the assignment intent for the mandated IST Coin: the persisted block reflects post-execution native balances, contract code,
and the token's mapping storage needed for correct \texttt{balanceOf} / allowance semantics after restart.

\section{Client Protocol and Byzantine-Tolerant Reads}
\subsection{Writes}
Clients broadcast signed transactions to all replicas and accept completion after receiving $f{+}1$ identical replies, rather than waiting for all replicas.
This prevents a Byzantine or slow replica from blocking progress.

\subsection{Read-only queries (balance)}
To support lightweight balance reads without heavy chain sync, we add a query request/response type.
Query replies include a \texttt{returnData} field (e.g., for ERC-20 \texttt{balanceOf}) and metadata about the responding replica head
(\texttt{headHeight}, \texttt{headHash}).
Clients accept a query result after receiving $f{+}1$ identical replies (tolerating Byzantine replies).
For read-after-write, clients may require replies from replicas whose head height is at least the decided index returned for the write.

\section{IST Coin Contract and Approval Frontrunning Mitigation}
The genesis deploys \texttt{ISTCoin}, an ERC-20 token with 2 decimals and a total supply of 100 million units.
To mitigate approval frontrunning, we enforce a guarded approval rule: a non-zero allowance must be reset to zero before being set to a new non-zero value.
This prevents the classic allowance-reduction race that allows withdrawing both old and new allowances.

\section{Threat Analysis and Countermeasures}
Table~\ref{tab:threats} summarizes key threats and defenses implemented across both stages.

\begin{table}[t]
\caption{Threats and protection mechanisms.}
\label{tab:threats}
\centering
\begin{tabularx}{\textwidth}{@{}lX X@{}}
\toprule
\textbf{Threat} & \textbf{Description} & \textbf{Protection} \\
\midrule
Message forgery / impersonation & Inject/alter consensus or client messages over UDP & Signatures; verify before deduplication (APL). \\
Replay / duplicate messages & Replaying old messages to confuse state & Per-sender message IDs; dedup after verification. \\
Byzantine leader proposing invalid payloads & Leader proposes non-verified client operations & Replicas validate proposals using verified pool; reject invalid blocks. \\
Byzantine clients forging transactions & Client submits tx spending others' funds & ECDSA signature verification; derive \texttt{from} from signature; enforce nonce. \\
Overspending / negative balances & Execute tx that makes balances negative & Pre-checks on balances; fee charging rules; abort semantics. \\
Approval frontrunning (ERC-20) & Allowance reduction race enables extra withdrawal & Guarded \texttt{approve}: require reset to zero before non-zero update. \\
Byzantine query replies & Replica lies on balance query & Client waits for $f{+}1$ identical query replies; include head metadata. \\
\bottomrule
\end{tabularx}
\end{table}

\section{Testing and Demonstration}
We use JUnit (see \texttt{README.md}) to demonstrate correctness and Byzantine tolerance:
\begin{itemize}
  \item Consensus safety/liveness under faults: \texttt{HotStuffIntegrationTest}, \texttt{ByzantineTest}, \texttt{ReplicaFailingTest}, \texttt{IntrusiveNetworkTest}.
  \item Byzantine clients/leaders and request validation: \texttt{ClientProtocolByzantineTest}, \texttt{ByzantineClientIntegrationTest}, \texttt{VerifiedBatchLeaderInjectionTest}.
  \item IST Coin: frontrunning-resistant approvals (\texttt{ISTCoinBesuExecutionTest}) and storage JSON round-trip (\texttt{ContractStoragePersistenceTest}).
  \item Gas-cap batching: \texttt{BlockBatcherTest}. Queries with $f{+}1$ identical replies: \texttt{ClientBalanceQueryIntegrationTest}, \texttt{ClientQueryByzantineReplyIntegrationTest},
  and read-after-write \texttt{ClientBalanceReadAfterWriteIntegrationTest}.
\end{itemize}
We also provide a multi-JVM demo script (\texttt{run-multijvm-demo.sh}).

\section{Dependability Guarantees}
\paragraph{Safety.}
With $n \ge 3f+1$ and at most $f$ Byzantine replicas, correct replicas decide the same sequence of blocks.
Transactions respect authorization (only owners can spend) and nonces prevent replays.
Account balances remain non-negative under the execution rules.

\paragraph{Liveness.}
Under partial synchrony, timeouts and view changes ensure progress: if the leader is faulty or slow, replicas move to the next view and a correct leader
eventually drives decisions.
Clients do not wait for all replicas; they accept completion on $f{+}1$ consistent replies.

\section{Limitations}
Our contract address derivation is project-specific (not Ethereum \texttt{CREATE}). Storage persistence is IST-focused and bounded; supporting arbitrary contracts
would require a more general storage tracing or full replay. We rely on clear documentation and reproducible tests/demos as part of the submission quality.

\section{Conclusion}
DepChain Stage~2 extends the Stage~1 BFT ordering layer with transaction semantics, an EVM execution environment, Ethereum-like fee accounting,
per-block gas limits, JSON persistence of accounts and IST Coin storage, and a hardened ERC-20 token deployed at genesis.
The system includes lightweight read queries with Byzantine-tolerant aggregation and tests for frontrunning mitigation, storage persistence, and BFT/client faults.

\bibliographystyle{splncs04}
\bibliography{refs}

\end{document}

